% Generated draft: requires mathematical and source review.
\documentclass[11pt]{article}
\usepackage[margin=25mm]{geometry}
\usepackage{fontspec}
\setmainfont{Latin Modern Roman}
\usepackage{amsmath,amssymb,mathtools,unicode-math}
\setmathfont{Latin Modern Math}
\usepackage{enumitem,needspace,xurl,hyperref}
\setlength{\parindent}{0pt}
\setlength{\parskip}{5pt}
\setlength{\emergencystretch}{3em}
\newcommand{\N}{\mathbb{N}}
\newcommand{\Z}{\mathbb{Z}}
\newcommand{\Q}{\mathbb{Q}}
\newcommand{\R}{\mathbb{R}}
\newcommand{\C}{\mathbb{C}}
\newcommand{\F}{\mathbb{F}}
\newcommand{\A}{\mathbb{A}}
\newcommand{\PP}{\mathbb{P}}
\newcommand{\E}{\mathbb{E}}
\newcommand{\eps}{\varepsilon}
\newcommand{\dd}{\,\mathrm{d}}
\begin{document}
\section*{UnsolvedMath ID: 30004059}
\textbf{Source problem code:} catalogue entry
\subsection{Definitions and mathematical statement}
The supplied record concerns a $G$-zip stack and the special fiber of a Shimura variety, but does not specify the underlying Shimura datum, group $G$, prime $p$, level $K$, or zip datum. The following formulation is therefore conditional on fixing the intended setup and the corresponding automorphic bundles.

Write $\mathcal Z$ for the intended $G$-zip stack and $S_K$ for the intended special fiber. Let $\Lambda$ be their common lattice of automorphic weights, and let $W\subseteq\Lambda$ be the set of weights for which the intended bundles are defined. For $\lambda\in W$, denote those bundles by $\mathcal V_{\mathcal Z}(\lambda)$ and $\mathcal V_K(\lambda)$, respectively. The weight sets described in the record are then
\[
C_{\mathrm{zip}}=\{\lambda\in W:H^0(\mathcal Z,\mathcal V_{\mathcal Z}(\lambda))\ne0\},\qquad
C_K(\mathbb F_p)=\{\lambda\in W:H^0(S_K,\mathcal V_K(\lambda))\ne0\}.
\]
Here the notation $C_K(\mathbb F_p)$ is retained exactly as supplied; the base fields and any intended base changes in these spaces of sections remain unspecified.

The supplied background describes the target as equality of saturated cones. With that interpretation, define, for $C\subseteq\Lambda$,
\[
\langle C\rangle=\{\lambda\in\Lambda:\text{there exists }n\in\Z_{>0}\text{ such that }n\lambda\in C\}.
\]
The question is whether, in the intended setup,
\[
\langle C_{\mathrm{zip}}\rangle=\langle C_K(\mathbb F_p)\rangle.
\]
Equivalently, is it true for every $\lambda\in\Lambda$ that
\[
\bigl(\exists n\in\Z_{>0}:n\lambda\in C_{\mathrm{zip}}\bigr)
\quad\Longleftrightarrow\quad
\bigl(\exists m\in\Z_{>0}:m\lambda\in C_K(\mathbb F_p)\bigr)?
\]
The weight lattice, admissible weights, bundle conventions, hypotheses on the geometric objects, and scope of quantification over Shimura-theoretic setups must be recovered from the original source before this conditional statement becomes a fully specified conjecture. The saturation convention above follows the supplied background; the record itself gives no explicit definition of the angle brackets.
\subsection{Short English statement}
Does a weight have a positive integer multiple admitting a nonzero automorphic section on the intended $G$-zip stack exactly when some positive integer multiple admits such a section on the intended Shimura special fiber? The geometric hypotheses, weight conventions, and base fields are missing from the supplied record, so this is a conditional statement requiring mathematical review.
\subsection{Sources}
\begin{itemize}
\item Supplied local UnsolvedMath registry record, permanent ID $30004059$, title \emph{Equality of Automorphic Weight Cones}, source problem code \texttt{catalogue entry}, in the Oberwolfach Reports Open Problems collection. The record lists proposed year $2019$ and identifies no individual proposer.
\item Supplied report citation: \emph{Arithmetic of Shimura Varieties (2020)}, Oberwolfach Reports, DOI \texttt{10.4171/owr/2019/2}, pp.\ 65--131. \url{https://doi.org/10.4171/owr/2019/2}. The supplied background attributes the problem to workshop open-problem discussions documented in this report.
\item Supplied literature source: Koskivirta, ``divisibility and the cone.'' \url{https://arxiv.org/abs/2211.16817}.
\item Supplied literature source: Irving--Koskivirta, ``automorphic forms on G-zips.'' \url{https://arxiv.org/abs/1810.05255}.
\item Literature assessment dated exactly \texttt{2026-08-22}: the supplied curation labels the problem partially solved and states that the saturated-cone equality remains conjectural in general, while inclusions, representation-theoretic descriptions, and several group or Shimura-variety cases are proved. This is supplied dated curation derived from the OWR record, not an independently verified literature assessment; time-sensitive status and cited claims require independent verification.
\end{itemize}
\end{document}
